
\documentclass{article}
\usepackage{colortbl}
\usepackage{makecell}
\usepackage{multirow}
\usepackage{supertabular}

\begin{document}

\newcounter{utterance}

\centering \large Interaction Transcript for game `cladder', experiment `full\_v1.5\_default', episode 3483 with SLED{-}BSU\_\_Qwen3.5{-}27B{-}sft{-}ep1\_\_prm{-}ep1{-}best\_of\_n.
\vspace{24pt}

{ \footnotesize  \setcounter{utterance}{1}
\setlength{\tabcolsep}{0pt}
\begin{supertabular}{c@{$\;$}|p{.15\linewidth}@{}p{.15\linewidth}p{.15\linewidth}p{.15\linewidth}p{.15\linewidth}p{.15\linewidth}}
    \# & \multicolumn{2}{c}{Player} && \multicolumn{2}{c}{Game Master} \\
    \hline

    \theutterance \stepcounter{utterance}  
    & & & \multicolumn{4}{p{0.6\linewidth}}{
        \cellcolor[rgb]{0.9,0.9,0.9}{
            \makecell[{{p{\linewidth}}}]{
                \texttt{\tiny{[P1$\langle$GM]}}
                \texttt{You are an expert in causal inference. The following question is not a typical commonsense query, but rather a meticulously designed question created by a professor specializing in causal inference, intended to assess the students' mastery of the course content.} \\
\texttt{Imagine a self{-}contained, hypothetical world with only the following conditions, and without any unmentioned factors or causal relationships: Tanning salon treatment has a direct effect on foot size. Going to the beach has a direct effect on foot size. The overall probability of tanning salon treatment is 10\%. For people not using tanning salon treatments, the probability of large feet is 13\%. For people who went to tanning salons, the probability of large feet is 70\%. Is large feet more likely than small feet overall?} \\
\texttt{Start your answer with 'Yes' or 'No', followed by additional reasoning or evidence to support your explanation.} \\
            }
        }
    }
     \\ \\

    \theutterance \stepcounter{utterance}  
    & \multicolumn{4}{p{0.6\linewidth}}{
        \cellcolor[rgb]{0.561,0.737,0.561}{
            \makecell[{{p{\linewidth}}}]{
                \texttt{\tiny{[P1$\rangle$GM]}}
                \texttt{Yes} \\
\\ 
\texttt{**Reasoning:**} \\
\\ 
\texttt{To determine whether large feet are more likely than small feet overall, we need to calculate the total probability of large feet across the entire population using the law of total probability.} \\
\\ 
\texttt{Let:} \\
\texttt{{-} \$T\$ = event of tanning salon treatment} \\
\texttt{{-} \$F\$ = event of large feet} \\
\\ 
\texttt{We are given:} \\
\texttt{{-} \$P(T) = 0.10\$ (10\% probability of tanning salon treatment)} \\
\texttt{{-} \$P(F \textbackslash{}mid \textbackslash{}neg T) = 0.13\$ (13\% probability of large feet without tanning salon)} \\
\texttt{{-} \$P(F \textbackslash{}mid T) = 0.70\$ (70\% probability of large feet with tanning salon)} \\
\\ 
\texttt{Using the law of total probability:} \\
\texttt{\$\$} \\
\texttt{P(F) = P(F \textbackslash{}mid T)P(T) + P(F \textbackslash{}mid \textbackslash{}neg T)P(\textbackslash{}neg T)} \\
\texttt{\$\$} \\
\\ 
\texttt{Substitute the values:} \\
\texttt{\$\$} \\
\texttt{P(F) = (0.70)(0.10) + (0.13)(0.90)} \\
\texttt{\$\$} \\
\texttt{\$\$} \\
\texttt{P(F) = 0.07 + 0.117 = 0.187} \\
\texttt{\$\$} \\
\\ 
\texttt{So the overall probability of large feet is 18.7\%.} \\
\\ 
\texttt{Since \$P(F) = 0.187 < 0.5\$, large feet are **not** more likely than small feet overall. Small feet have probability \$1 {-} 0.187 = 0.813\$, which is much higher.} \\
\\ 
\texttt{Therefore, the correct answer is **No**.} \\
            }
        }
    }
    & & \\ \\

    \theutterance \stepcounter{utterance}  
    & & & \multicolumn{2}{p{0.3\linewidth}}{
        \cellcolor[rgb]{0.9,0.9,0.9}{
            \makecell[{{p{\linewidth}}}]{
                \texttt{\tiny{[GM$|$GM]}}
                \texttt{Yes} \\
            }
        }
    }
    & & \\ \\

    \theutterance \stepcounter{utterance}  
    & & & \multicolumn{2}{p{0.3\linewidth}}{
        \cellcolor[rgb]{0.9,0.9,0.9}{
            \makecell[{{p{\linewidth}}}]{
                \texttt{\tiny{[GM$|$GM]}}
                \texttt{no} \\
            }
        }
    }
    & & \\ \\

    \theutterance \stepcounter{utterance}  
    & & & \multicolumn{2}{p{0.3\linewidth}}{
        \cellcolor[rgb]{0.9,0.9,0.9}{
            \makecell[{{p{\linewidth}}}]{
                \texttt{\tiny{[GM$|$GM]}}
                \texttt{game\_result = LOSE} \\
            }
        }
    }
    & & \\ \\

\end{supertabular}
}

\end{document}
